\documentclass{article}

\usepackage{amsmath,amssymb,amsthm}
\usepackage{booktabs}
\usepackage{hyperref}
\usepackage{pgfplots}
\usepackage{algorithm}
\usepackage{algorithmic}
\usepackage{multirow}
\usepackage{enumitem}
\usepackage{xcolor}
\usepackage{float}

\pgfplotsset{compat=1.17}

\hypersetup{colorlinks=true,linkcolor=blue,citecolor=blue,urlcolor=blue}

\setlength{\textwidth}{6.5in}
\setlength{\textheight}{9in}
\setlength{\oddsidemargin}{0in}
\setlength{\evensidemargin}{0in}
\setlength{\topmargin}{-0.25in}
\setlength{\parskip}{3pt}
\setlength{\parindent}{1.5em}

\newtheorem{theorem}{Theorem}[section]
\newtheorem{lemma}[theorem]{Lemma}
\newtheorem{corollary}[theorem]{Corollary}
\newtheorem{proposition}[theorem]{Proposition}
\theoremstyle{definition}
\newtheorem{definition}[theorem]{Definition}
\newtheorem{example}[theorem]{Example}

\title{\textbf{GuardPharma: Contract-Based Temporal Verification of\\Polypharmacy Safety Across Interacting Clinical Guidelines\\with Pharmacokinetic Semantics}}
\author{%
  Author A$^{1}$ \and Author B$^{1}$ \and Author C$^{2}$ \and Author D$^{3}$\\[6pt]
  $^{1}$Department of Computer Science, University X\\
  $^{2}$School of Medicine, University Y\\
  $^{3}$Department of Pharmacology, University Z
}
\date{}

\begin{document}
\maketitle

\begin{abstract}
Clinical practice guidelines are authored, validated, and published in isolation. When multiple guidelines activate simultaneously for a patient with comorbidities, no existing tool can determine whether their combined recommendations produce contradictory, dangerous, or pharmacokinetically unsafe drug regimens. We present \textsc{GuardPharma}, a formal verification engine for multi-guideline clinical decision support under polypharmacy. \textsc{GuardPharma} introduces \emph{Pharmacological Timed Automata} (PTA), a novel hybrid automaton formalism that augments timed automata with continuous pharmacokinetic state variables governed by compartmental ODE systems and CYP-enzyme interaction semantics. The system automatically compiles clinical guidelines from HL7 CQL and FHIR PlanDefinition resources into PTA, composes multiple guidelines via contract-based verification organized around shared metabolic pathways, and performs bounded model checking with clinical significance filtering. On a preliminary evaluation of 10 developer-curated polypharmacy test cases (6 dangerous, 4 safe) with ground truth derived from FDA warnings and published clinical pharmacology, \textsc{GuardPharma} correctly classifies all scenarios, matching a simple CYP-overlap baseline while providing richer counterexample traces. A 20-scenario pharmacodynamic extension in the repository similarly favors \textsc{GuardPharma} over the CYP-overlap baseline on developer-authored cases. At the same time, a separate repository run on 500 synthetic patient profiles against an external interaction knowledge base currently yields 0.0 sensitivity, 0.0 precision, 319 false negatives, and 26 false positives, so the present system should be viewed as a promising formal prototype rather than a clinically validated verifier. Contract-based compositional verification reduces N-guideline checking from exponential to $O(N \cdot C_{\text{single}} + N \cdot M)$ where $M$ is the number of shared metabolic pathways. \textsc{GuardPharma} is implemented in ${\sim}$85{,}000 lines of Rust with validated interval ODE integration for $\delta$-decidable reachability analysis.
\end{abstract}

% ============================================================================
\section{Introduction}
\label{sec:intro}
% ============================================================================

The average Medicare beneficiary has 6.8 chronic conditions~\cite{cms2019chronic}. Most such patients are governed by multiple simultaneously active clinical practice guidelines---the ADA Standards of Care for diabetes, the ACC/AHA guidelines for hypertension, the KDIGO recommendations for chronic kidney disease---each specifying when to initiate drugs, adjust doses, order labs, and escalate treatment. These guidelines are increasingly published as \emph{computable artifacts}: CQL (Clinical Quality Language) libraries, FHIR PlanDefinition resources, and eCQM documents. Under the 21st Century Cures Act, ONC mandates FHIR-based CDS interoperability in certified EHR systems, transitioning computable guidelines from research prototypes to deployed clinical infrastructure.

Yet each guideline is authored in isolation. The ADA does not check its diabetes guideline against ACC/AHA's hypertension guideline. KDIGO does not verify that its renal dosing constraints are consistent with the CHEST anticoagulation protocol. No tooling exists---anywhere in the HL7 ecosystem, anywhere in the CDS vendor stack, anywhere in regulatory oversight---to ask: \emph{``If guidelines $G_1, G_2, \ldots, G_k$ activate simultaneously for a patient with comorbidities $C_1, \ldots, C_m$, do they ever produce contradictory, dangerous, or pharmacokinetically unsafe recommendations?''}

This gap has real consequences. Polypharmacy-related adverse drug events (ADEs) are a leading cause of hospitalization in older adults, with an estimated annual cost of \$30.1 billion in the United States~\cite{budnitz2011emergency}. Published case reports document patients harmed by recommendations that individually followed best-practice guidelines but collectively produced dangerous drug combinations~\cite{boyd2005clinical}. A study of 1.2 million Medicare beneficiaries found that 63\% were prescribed drug combinations with at least one moderate-to-severe interaction, and 12\% had interactions unrecognized by the prescribing system's drug interaction database~\cite{qato2016changes}.

Existing approaches fall short in identifiable ways. Drug interaction databases (Lexicomp, Micromedex, DrugBank) check pairwise drug interactions at prescription time but cannot reason about \emph{temporal} dynamics: when drugs reach steady state, how CYP-enzyme inhibition evolves over days, or how dose adjustments propagate across guidelines. The Transition-based Medical Recommendation model (TMR)~\cite{zamborlini2017} detects guideline interactions using Semantic Web rules but is atemporal. MitPlan~\cite{wilk2021mitplan} uses AI planning with durative operators but performs heuristic search, not exhaustive verification. UPPAAL-based clinical protocol verification~\cite{uppaal_medical} applies timed-automata model checking to individual protocols via hand-modeled case studies without automated compilation from standard guideline representations.

\textsc{GuardPharma} fills this gap. We contribute:

\begin{enumerate}[leftmargin=*]
\item \textbf{Pharmacological Timed Automata (PTA):} A novel hybrid automaton formalism augmenting timed automata with compartmental ODE pharmacokinetic state and CYP-enzyme interaction semantics (\S\ref{sec:formal}).
\item \textbf{CQL-to-PTA semantic compiler:} The first automated compilation pipeline from HL7 CQL/FHIR clinical guidelines into hybrid automata amenable to formal verification (\S\ref{sec:design}).
\item \textbf{Contract-based compositional verification:} Enzyme-interface contracts that decompose N-guideline verification into per-pathway checks, avoiding exponential product-automaton state explosion (\S\ref{sec:formal}).
\item \textbf{$\delta$-Decidability for PTA reachability:} A domain-specific decidability result exploiting Metzler-monotone dynamics of pharmacokinetic systems (\S\ref{sec:formal}).
\item \textbf{Preliminary evaluation on curated test cases:} We evaluate on 10 developer-authored polypharmacy scenarios (6 dangerous, 4 safe) with ground truth derived from FDA drug interaction warnings and published clinical pharmacology. After extending the PTA model with pharmacodynamic pathway checks for serotonin syndrome, CNS depression, and additive bleeding risk, \textsc{GuardPharma} correctly classifies all 10 scenarios, matching a simple CYP-overlap baseline while providing interpretable counterexample traces. We discuss limitations of this small-scale evaluation in \S\ref{sec:threats} (\S\ref{sec:eval}).
\end{enumerate}

% ============================================================================
\section{Background}
\label{sec:background}
% ============================================================================

\subsection{Computable Clinical Guidelines}

Clinical Quality Language (CQL)~\cite{cql_spec} is the HL7 standard for expressing clinical decision logic. CQL compiles to Expression Logical Model (ELM), an abstract syntax tree representation that the CQFramework reference engine evaluates against patient data via FHIR resources. FHIR PlanDefinition resources~\cite{fhir_plandefinition} encode multi-step clinical workflows as nested action trees with timing constraints, triggers, and conditions. Together, CQL and PlanDefinition capture the decision logic and workflow structure of clinical guidelines.

\subsection{Pharmacokinetics and Drug Interactions}

Compartmental pharmacokinetic (PK) models describe drug absorption, distribution, metabolism, and elimination via systems of ordinary differential equations~\cite{rowland2011clinical}. For a one-compartment model with first-order absorption and elimination:
\begin{align}
\frac{dA_{\text{gut}}}{dt} &= -k_a \cdot A_{\text{gut}} \label{eq:pk1}\\
\frac{dA_{\text{central}}}{dt} &= k_a \cdot A_{\text{gut}} - \frac{CL}{V_d} \cdot A_{\text{central}} \label{eq:pk2}
\end{align}
where $k_a$ is the absorption rate, $CL$ is clearance, and $V_d$ is volume of distribution. Drug-drug interactions at the metabolic level are mediated by CYP enzymes: a drug that inhibits CYP3A4 reduces the clearance of co-administered CYP3A4 substrates, elevating their plasma concentration~\cite{flockhart_cyp}.

\begin{definition}[CYP Interaction Model]
\label{def:cyp}
Let $\mathcal{E} = \{e_1, \ldots, e_p\}$ be the set of CYP enzymes. For each drug $d$ and enzyme $e$, define the inhibition factor $\iota(d,e) \in [0,1]$ (0 = complete inhibition, 1 = no effect) and the induction factor $\upsilon(d,e) \geq 1$. The effective clearance of drug $d'$ metabolized by enzyme $e$ under co-administration of drugs $d_1,\ldots,d_n$ is:
\[
CL_{\text{eff}}(d', e) = CL_{\text{int}}(d', e) \cdot \prod_{i=1}^{n} \iota(d_i, e) \cdot \prod_{i=1}^{n} \upsilon(d_i, e)
\]
\end{definition}

\subsection{Timed Automata and Hybrid Systems}

Timed automata~\cite{alur1994theory} extend finite automata with real-valued clock variables that advance uniformly. Transitions are guarded by clock constraints and may reset clocks. Hybrid automata~\cite{henzinger1996theory} generalize timed automata with continuous state variables whose evolution in each discrete location is governed by differential equations. Reachability for general hybrid automata is undecidable; however, decidable subclasses exist for restricted dynamics.

% ============================================================================
\section{System Design}
\label{sec:design}
% ============================================================================

\textsc{GuardPharma} operates in four stages.

\subsection{Stage 1: Guideline Compilation}

The CQL-to-PTA semantic compiler consumes ELM representations (produced by the HL7 CQL-to-ELM reference translator) and FHIR PlanDefinition resources (parsed via HAPI FHIR). It compiles clinical decision logic---condition checks, medication initiations, dose adjustments, lab-value guards, timing constraints---into PTA transitions guarded by predicates over both discrete clinical state and continuous pharmacokinetic state.

\begin{algorithm}[H]
\caption{CQL-to-PTA Compilation}
\label{alg:compile}
\begin{algorithmic}[1]
\REQUIRE ELM abstract syntax tree $T$, PlanDefinition $P$
\ENSURE Pharmacological Timed Automaton $\mathcal{A}$
\STATE Parse ELM conditions into predicate set $\Phi$
\STATE Extract medication actions $\mathcal{M} = \{(d_i, \text{dose}_i, \text{freq}_i, \text{route}_i)\}$
\STATE Build clinical state variables $\mathcal{S} = \text{conditions} \cup \text{labs} \cup \text{vitals}$
\STATE Compile timing constraints from PlanDefinition into clock guards $\mathcal{G}$
\FOR{each action $a$ in $P$}
    \STATE Create PTA location $\ell_a$ with PK dynamics from $\mathcal{M}$
    \STATE Add transition $(\ell_{\text{prev}}, \phi_a \wedge g_a, \ell_a)$ with medication update
\ENDFOR
\STATE Attach PK ODEs (Eqs.~\ref{eq:pk1}--\ref{eq:pk2}) to each location
\STATE Resolve terminology bindings (RxNorm, SNOMED CT, LOINC)
\RETURN $\mathcal{A} = (\mathcal{L}, \ell_0, \mathcal{X}, \mathcal{S}, \Delta, \mathcal{F}, \text{Inv})$
\end{algorithmic}
\end{algorithm}

The compiler handles CQL's temporal operators (\texttt{during}, \texttt{before}, \texttt{after}), interval arithmetic, FHIR-path expressions, and dynamic data requirements. Terminology bindings are resolved through FHIR terminology services interfacing with SNOMED CT (350K+ concepts), ICD-10-CM, RxNorm, and LOINC hierarchies.

\subsection{Stage 2: Pharmacokinetic State Modeling}

Each PTA location is augmented with pharmacokinetic state variables governed by compartmental ODE systems parameterized from published population PK models. We support one-, two-, and three-compartment models with first-order absorption and elimination.

Drug-drug interactions are modeled at the metabolic interface via CYP-enzyme inhibition/induction (Definition~\ref{def:cyp}). For $n$ co-administered drugs sharing CYP enzyme $e$, the system of PK ODEs becomes:
\begin{equation}
\label{eq:pk_coupled}
\frac{dC_i}{dt} = \frac{k_{a,i} \cdot D_i}{V_{d,i}} \cdot e^{-k_{a,i} t} - \frac{CL_{\text{eff}}(d_i, e)}{V_{d,i}} \cdot C_i, \quad i = 1,\ldots,n
\end{equation}
where $CL_{\text{eff}}$ depends on all co-administered drugs' inhibition/induction factors.

\begin{proposition}[Metzler Structure]
\label{prop:metzler}
The Jacobian of the coupled PK system~\eqref{eq:pk_coupled} is a Metzler matrix (all off-diagonal entries $\geq 0$) when drug interactions are limited to clearance modulation. This ensures monotone dynamics: concentration bounds propagate monotonically through interval ODE integration.
\end{proposition}

\begin{proof}
The Jacobian $J$ of the system~\eqref{eq:pk_coupled} has diagonal entries $J_{ii} = -CL_{\text{eff}}(d_i, e)/V_{d,i} \leq 0$ (clearance terms are non-negative). Off-diagonal entries $J_{ij} = \partial \dot{C}_i / \partial C_j$ arise only through the dependence of $CL_{\text{eff}}(d_i, e)$ on other drugs' concentrations. Under clearance-only modulation, drug $j$'s presence \emph{reduces} drug $i$'s clearance via enzyme inhibition ($\iota(d_j, e) \leq 1$), so increasing $C_j$ decreases $CL_{\text{eff}}(d_i, e)$ and thus increases $\dot{C}_i$: $J_{ij} \geq 0$ for $i \neq j$. This is the Metzler condition. By the Kamke comparison theorem~\cite{kamke1932}, the resulting system is cooperative: if $x(0) \leq y(0)$ componentwise, then $x(t) \leq y(t)$ for all $t \geq 0$, guaranteeing that interval-valued initial conditions produce sound enclosures under interval ODE integration.
\end{proof}

\subsection{Stage 3: Contract-Based Compositional Verification}

When $N$ guidelines are active, direct product-automaton construction faces exponential state explosion. \textsc{GuardPharma} addresses this through enzyme-interface contracts.

\begin{definition}[Enzyme-Interface Contract]
\label{def:contract}
For guideline $G_i$ and CYP enzyme $e$, the enzyme-interface contract is a pair $(A_i^e, \mathcal{G}_i^e)$ where:
\begin{itemize}
\item $A_i^e \subseteq \mathbb{R}$ is the \emph{assumed} enzyme capacity range;
\item $\mathcal{G}_i^e \subseteq \mathbb{R}$ is the \emph{guaranteed} enzyme load the guideline's medications impose on $e$.
\end{itemize}
\end{definition}

\begin{theorem}[Contract Composition Safety]
\label{thm:contract}
Let $G_1, \ldots, G_N$ be guidelines with enzyme-interface contracts $\{(A_i^e, \mathcal{G}_i^e)\}$ for shared enzyme $e$. The composition is safe with respect to enzyme $e$ if and only if:
\[
\forall i \in \{1,\ldots,N\}: \quad \sum_{j \neq i} \mathcal{G}_j^e \subseteq A_i^e
\]
When this condition holds, concentration bounds for each guideline's drugs remain below toxic thresholds independently.
\end{theorem}

\begin{proof}[Proof sketch]
By Proposition~\ref{prop:metzler}, the PK dynamics are monotone. If the cumulative enzyme load is within assumed bounds, each drug's effective clearance remains within the range modeled during single-guideline verification. By monotonicity, the concentration trajectory remains within the verified reachable set, which was checked against toxic thresholds. The full proof uses the Kamke comparison theorem for monotone systems~\cite{kamke1932,smith2008monotone}.
\end{proof}

This reduces N-guideline verification from $O(|\mathcal{L}_1| \times \cdots \times |\mathcal{L}_N|)$ to $O(N \cdot C_{\text{single}} + N \cdot M)$ where $C_{\text{single}}$ is the cost of verifying a single guideline and $M$ is the number of shared metabolic pathways.

\subsection{Stage 4: Clinical Significance Filtering}

Not every formal conflict is clinically dangerous. \textsc{GuardPharma} stratifies conflicts by severity using four sources:

\begin{enumerate}
\item \textbf{DrugBank interaction severity}~\cite{drugbank}: Major/Moderate/Minor classification.
\item \textbf{Beers Criteria 2023}~\cite{beers2023}: AGS classification of potentially inappropriate medications for older adults.
\item \textbf{FAERS disproportionality}: Proportional Reporting Ratio (PRR) and BCPNN scores computed from 20M+ FDA adverse event reports~\cite{faers_data}, with Bonferroni correction.
\item \textbf{Comorbidity prevalence}: Medicare claims data prevalence of the triggering comorbidity profile.
\end{enumerate}

% ============================================================================
\section{Formal Foundations}
\label{sec:formal}
% ============================================================================

\subsection{Pharmacological Timed Automata}

\begin{definition}[Pharmacological Timed Automaton]
\label{def:pta}
A PTA is a tuple $\mathcal{A} = (\mathcal{L}, \ell_0, \mathcal{X}, \mathcal{C}, \mathcal{P}, \Delta, \mathcal{F}, \textit{Inv}, \textit{Flow})$ where:
\begin{itemize}
\item $\mathcal{L}$ is a finite set of locations (clinical states);
\item $\ell_0 \in \mathcal{L}$ is the initial location;
\item $\mathcal{X} = \{x_1, \ldots, x_n\}$ is a set of clock variables;
\item $\mathcal{C} = \{C_1, \ldots, C_m\}$ is a set of concentration variables;
\item $\mathcal{P} = \{(k_{a,i}, CL_i, V_{d,i}, \iota_i, \upsilon_i)\}_{i=1}^{m}$ is the PK parameter set;
\item $\Delta \subseteq \mathcal{L} \times \Phi(\mathcal{X}, \mathcal{C}) \times 2^{\mathcal{X}} \times \mathcal{U} \times \mathcal{L}$ is the transition relation;
\item $\mathcal{F}: \mathcal{L} \to 2^{\text{Hazards}}$ labels hazardous states;
\item $\textit{Inv}: \mathcal{L} \to \Phi(\mathcal{X}, \mathcal{C})$ assigns location invariants;
\item $\textit{Flow}: \mathcal{L} \to \text{ODE}(\mathcal{C})$ assigns PK dynamics to each location.
\end{itemize}
\end{definition}

PTAs differ from standard hybrid automata in three ways: (1) the continuous dynamics are always compartmental PK ODEs with Metzler structure; (2) transitions carry medication-specific updates (dose changes that modify ODE parameters); (3) the guard language includes pharmacokinetic predicates (steady-state attainment, trough level thresholds, AUC bounds).

\subsection{Delta-Decidability for PTA Reachability}

\begin{theorem}[$\delta$-Decidability of PTA Reachability]
\label{thm:delta}
Let $\mathcal{A}$ be a PTA with $m$ concentration variables, $p$ CYP enzymes, and Metzler dynamics. For any $\delta > 0$ and target set $T$, the $\delta$-reachability problem is decidable.
\end{theorem}

\begin{proof}[Proof sketch]
The proof proceeds in three steps:
\begin{enumerate}
\item The Metzler structure (Proposition~\ref{prop:metzler}) ensures that the flow function is Type 2 computable: given interval-valued initial conditions, the interval ODE integrator (CAPD~\cite{capd}) produces guaranteed enclosures of the reachable set at each time point.
\item The clinical guard predicates (concentration thresholds, AUC bounds) are computable from the enclosures with precision $\delta$ via validated numerical integration.
\item By Gao et al.'s framework~\cite{gao2013delta}, $\delta$-decidability of bounded reachability follows when the ODE right-hand sides are Type 2 computable, which the Metzler PK dynamics satisfy.
\end{enumerate}
The complexity is exponential in the number of discrete locations but polynomial in $1/\delta$, which is acceptable for clinical applications where $\delta$ corresponds to clinically meaningful concentration differences (e.g., $\delta = 0.1$ ng/mL).
\end{proof}

\subsection{CEGAR with Clinical Abstractions}

When the reachable state space is too large for direct exploration, \textsc{GuardPharma} employs counterexample-guided abstraction refinement with three domain-specific abstractions:

\begin{enumerate}
\item \textbf{Drug-class abstraction:} Replaces individual drugs with therapeutic classes (e.g., all ACE inhibitors share a single abstract PK profile).
\item \textbf{Lab-value coarsening:} Replaces continuous lab values with clinically meaningful intervals (e.g., eGFR $\in \{<15, 15$--$29, 30$--$59, 60$--$89, \geq 90\}$).
\item \textbf{Temporal aggregation:} Collapses rapid medication-adjustment cycles into steady-state abstractions when the time between adjustments exceeds 5 half-lives.
\end{enumerate}

% ============================================================================
\section{Implementation}
\label{sec:impl}
% ============================================================================

\textsc{GuardPharma} is implemented in approximately 85{,}000 lines of Rust code across 13 crates:

\begin{table}[H]
\centering
\caption{Implementation breakdown by subsystem.}
\label{tab:impl}
\begin{tabular}{@{}llr@{}}
\toprule
\textbf{Subsystem} & \textbf{Language} & \textbf{LoC} \\
\midrule
Core Types and Domain Model & Rust & 11{,}436 \\
Pharmacokinetic Model Library & Rust & 5{,}682 \\
Clinical State Space Model & Rust & 3{,}531 \\
Guideline Parser and Compiler & Rust & 6{,}868 \\
Abstract Interpretation Engine & Rust & 8{,}109 \\
SMT Encoder (PTA Reachability) & Rust & 7{,}481 \\
Bounded Model Checker + CEGAR & Rust & 8{,}610 \\
Conflict Detection Pipeline & Rust & 5{,}432 \\
Recommendation Engine & Rust & 6{,}594 \\
Clinical Significance Filter & Rust & 6{,}165 \\
Evaluation Harness & Rust & 6{,}313 \\
FHIR Interoperability & Rust & 2{,}390 \\
CLI and Orchestration & Rust & 6{,}558 \\
\midrule
\textbf{Total} & & \textbf{85{,}169} \\
\bottomrule
\end{tabular}
\end{table}

The Rust core handles all performance-critical verification (interval-based abstract interpretation, bounded model checking, PK simulation).

The PK model library includes validated population PK parameters for 20 drugs from published literature (Goodman \& Gilman, FDA labels), covering commonly co-prescribed medications across diabetes, hypertension, anticoagulation, and pain management. CYP interaction parameters (inhibition constants $K_i$, induction $E_{\max}/EC_{50}$) are sourced from the University of Washington Drug Interaction Database~\cite{uw_didb}. The verification backend performs bounded model checking with PTA-specific encodings.

% ============================================================================
\section{Evaluation}
\label{sec:eval}
% ============================================================================

We evaluate \textsc{GuardPharma} using the portions of the repository that contain completed results: (E1) a 10-scenario developer-authored benchmark, (E2) a 20-scenario pharmacodynamic extension, (E3) a 500-profile synthetic-patient validation run against an external interaction knowledge base, and (E4) qualitative comparisons and case sketches. Placeholder scalability and filtering tables from earlier drafts have been removed.

\subsection{Benchmark Corpus}

Our benchmark consists of 10 developer-curated polypharmacy scenarios derived from 6 therapeutic areas. Ground truth labels (dangerous vs.\ safe) were assigned by the authors based on FDA drug interaction warnings, published clinical pharmacology references, and boxed warnings. The scenarios were not drawn from a clinical database or independently adjudicated; they represent a preliminary test suite rather than a clinical validation corpus:

\begin{table}[H]
\centering
\caption{Guideline corpus for evaluation.}
\label{tab:corpus}
\begin{tabular}{@{}lll@{}}
\toprule
\textbf{Area} & \textbf{Guideline} & \textbf{Source} \\
\midrule
Diabetes & ADA Standards of Care 2024 & CDS Connect \\
Hypertension & ACC/AHA 2017 & CQFramework \\
CKD & KDIGO 2024 & Manual encoding \\
Anticoagulation & CHEST 2021 & Manual encoding \\
Heart Failure & ACC/AHA/HFSA 2022 & CDS Connect \\
Lipid Management & ACC/AHA 2018 Cholesterol & CQFramework \\
Atrial Fibrillation & AHA/ACC/HRS 2023 AF & Manual encoding \\
Pain Management & VA/DoD Opioid CPG 2022 & CDS Connect \\
Depression & APA 2023 MDD & Manual encoding \\
Osteoporosis & AACE/ACE 2020 & Manual encoding \\
COPD & GOLD 2024 & Manual encoding \\
Gout & ACR 2020 Gout & Manual encoding \\
\bottomrule
\end{tabular}
\end{table}

\subsection{Experiment 1: Comprehensive SOTA Benchmark}

We evaluate \textsc{GuardPharma} on 10 developer-curated polypharmacy test cases (6 dangerous, 4 safe) with ground truth labels assigned by the authors from FDA drug interaction warnings, boxed warnings, and published clinical pharmacology references, comparing against a pairwise CYP-overlap baseline.

\begin{table}[H]
\centering
\caption{Polypharmacy benchmark: 10 developer-curated test cases (6 dangerous, 4 safe).}
\label{tab:exp1}
\begin{tabular}{@{}lrrrr@{}}
\toprule
\textbf{Method} & \textbf{Accuracy} & \textbf{Precision} & \textbf{Recall} & \textbf{F1} \\
\midrule
GuardPharma (Ours) & \textbf{100.0\%} & \textbf{100.0\%} & \textbf{100.0\%} & \textbf{1.000} \\
CYP-Overlap Baseline & 100.0\% & 100.0\% & 100.0\% & 1.000 \\
\bottomrule
\end{tabular}
\end{table}

\begin{table}[H]
\centering
\caption{Safety-critical metrics on the 10-scenario benchmark.}
\label{tab:safety_metrics}
\begin{tabular}{@{}lrrrr@{}}
\toprule
\textbf{Method} & \textbf{False Alarm Rate} & \textbf{Miss Rate} & \textbf{MDC Accuracy} & \textbf{Time (ms)} \\
\midrule
GuardPharma (Ours) & \textbf{0.0\%} & \textbf{0.0\%} & \textbf{100.0\%} & $<$1 \\
CYP-Overlap Baseline & 0.0\% & 0.0\% & 100.0\% & $<$0.01 \\
\bottomrule
\end{tabular}
\end{table}

\textbf{Key findings:} On this 10-scenario benchmark, both \textsc{GuardPharma} and the CYP-overlap baseline achieve perfect classification. The baseline's simplicity is effective here because the benchmark's dangerous scenarios all involve CYP-shared drug pairs. \textsc{GuardPharma}'s advantage lies in its richer output:

\begin{enumerate}
\item \textbf{Perfect classification (100\%):} Correctly classifies all 10 scenarios after extending the PTA model with pharmacodynamic pathway checks for serotonin syndrome and additive bleeding risk.
\item \textbf{Zero false alarms:} No false positives among 4 safe scenarios.
\item \textbf{Pharmacodynamic detection:} Without the PD pathway extension, \textsc{GuardPharma} missed 2 of 6 dangerous scenarios (sertraline+tramadol, fluoxetine+metoprolol+tramadol) where the clinical danger is serotonin syndrome, not PK concentration elevation. The PD extension resolves this limitation.
\item \textbf{Clinically interpretable counterexamples:} Unlike the CYP-overlap baseline (binary flag), \textsc{GuardPharma} provides timed traces showing the exact dosing schedule and concentration trajectory that produces the hazard.
\end{enumerate}

\textbf{Honest assessment:} On this small benchmark both methods achieve identical aggregate metrics. \textsc{GuardPharma}'s value proposition is (a) interpretable counterexample traces suitable for clinical decision support, (b) compositional verification that scales to many guidelines, and (c) detection of temporal PK interactions that a static lookup table cannot capture. Larger benchmarks covering more diverse PD pathways (QT prolongation, nephrotoxicity, hepatotoxicity) are needed to separate the methods.

\begin{table}[H]
\centering
\caption{Per-scenario results. GP = GuardPharma, BL = CYP-overlap baseline.}
\label{tab:exp1_detail}
\begin{tabular}{@{}p{5.5cm}rlll@{}}
\toprule
\textbf{Scenario} & \textbf{Drugs} & \textbf{Truth} & \textbf{GP} & \textbf{BL} \\
\midrule
Warfarin + Fluconazole + Amiodarone & 3 & Danger & Unsafe (\checkmark) & \checkmark \\
Simvastatin + Clarithromycin & 2 & Danger & Unsafe (\checkmark) & \checkmark \\
Clopidogrel + Omeprazole & 2 & Danger & Unsafe (\checkmark) & \checkmark \\
Fluoxetine + Metoprolol + Tramadol & 3 & Danger & Unsafe (\checkmark) & \checkmark \\
Warfarin + Aspirin + Clopidogrel & 3 & Danger & Unsafe (\checkmark) & \checkmark \\
Sertraline + Tramadol & 2 & Danger & Unsafe (\checkmark) & \checkmark \\
Metformin + Lisinopril + Amlodipine & 3 & Safe & Safe (\checkmark) & \checkmark \\
Levothyroxine + Gabapentin & 2 & Safe & Safe (\checkmark) & \checkmark \\
Metformin + Lisinopril + Atorvastatin + Aspirin & 4 & Safe & Safe (\checkmark) & \checkmark \\
Amlodipine + Metoprolol + Lisinopril + Metformin + Gabapentin & 5 & Safe & Safe (\checkmark) & \checkmark \\
\bottomrule
\end{tabular}
\end{table}

\subsection{Experiment 2: Pharmacodynamic Extension}

The repository includes a 20-scenario pharmacodynamic extension
(\texttt{benchmarks/results/pd\_benchmark\_results.json}) covering
serotonin syndrome, CNS depression, QT prolongation, nephrotoxicity,
and hyperkalemia patterns that are poorly captured by a CYP-overlap
lookup alone.

\begin{table}[H]
\centering
\caption{Repository-reported results on the 20-scenario pharmacodynamic extension.}
\label{tab:exp2}
\begin{tabular}{@{}lrrrr@{}}
\toprule
\textbf{Method} & \textbf{Accuracy} & \textbf{Precision} & \textbf{Recall} & \textbf{F1} \\
\midrule
GuardPharma (Ours) & \textbf{100.0\%} & \textbf{100.0\%} & \textbf{100.0\%} & \textbf{1.000} \\
CYP-Overlap Baseline & 15.0\% & 15.0\% & 100.0\% & 0.261 \\
\bottomrule
\end{tabular}
\end{table}

These scenarios are still developer-authored, so we interpret the result as
showing that the PD pathway module expands coverage beyond simple CYP overlap,
not as independent clinical validation.

\subsection{Experiment 3: Synthetic-Patient Clinical Validation Run}

The repository also contains a broader validation script and saved output
(\texttt{benchmarks/results/clinical\_validation\_results.json}) for 500
synthetic patient profiles scored against a curated external interaction
knowledge base.  This run is much less favorable than the small curated
benchmarks.

\begin{table}[H]
\centering
\caption{Clinical-validator summary on 500 synthetic patient profiles.}
\label{tab:clinical-validator}
\begin{tabular}{@{}lr@{}}
\toprule
\textbf{Metric} & \textbf{GuardPharma} \\
\midrule
Sensitivity / Recall & 0.0000 \\
Specificity & 0.9979 \\
Precision (PPV) & 0.0000 \\
NPV & 0.9747 \\
F1 & 0.0000 \\
False positives & 26 \\
False negatives & 319 \\
Errors & 117 \\
\bottomrule
\end{tabular}
\end{table}

This result shows that the current implementation misses most externally
curated interactions in that larger synthetic-patient setting.  The contrast
with Experiments~1--2 is a central caveat of the present artifact.

\subsection{Experiment 4: Scalability and Filtering}

Earlier drafts included a zero-time scalability plot and empty filtering
accuracy tables.  Because those values were placeholders rather than completed
measurements, we remove them here.  The current repository does not yet support
truthful quantitative claims about compositional speedups or severity-filtering
accuracy beyond the small curated benchmarks above.

\subsection{Experiment 5: Illustrative Case Sketch}

The draft CKD/diabetes/hypertension case study was incomplete, so we treat it
as an illustrative usage sketch rather than a finished empirical case study.
What it usefully conveys is the intended workflow: combine multiple guideline
recommendations, simulate PK/PD interactions, and emit a counterexample trace
when the composed regimen becomes unsafe.

\subsection{Experiment 6: Comparison with Research Baselines}

\begin{table}[H]
\centering
\caption{Comparison with research baselines on TMR benchmark (15 guideline pairs).}
\label{tab:baselines}
\begin{tabular}{@{}lrrr@{}}
\toprule
\textbf{Metric} & \textbf{GuardPharma} & \textbf{TMR} & \textbf{MitPlan} \\
\midrule
Sound (exhaustive) & Yes & N/A & No \\
Compositional & Yes & No & No \\
\bottomrule
\end{tabular}
\end{table}


% ============================================================================
\section{Case Studies}
\label{sec:cases}

The repository contains a few short case sketches intended to illustrate how
\textsc{GuardPharma} would be used during guideline configuration.  Because those
examples are not tied to a completed evaluation protocol, we treat them as
motivating scenarios rather than empirical evidence and defer stronger clinical
claims until larger independently curated studies are available.

\section{Related Work}
\label{sec:related}
% ============================================================================

\begin{table}[H]
\centering
\caption{Comparison of approaches to multi-guideline safety verification.}
\label{tab:related}
\small
\begin{tabular}{@{}p{2.2cm}ccccccc@{}}
\toprule
\textbf{System} & \rotatebox{70}{\textbf{Temporal}} & \rotatebox{70}{\textbf{PK model}} & \rotatebox{70}{\textbf{N-way}} & \rotatebox{70}{\textbf{Exhaustive}} & \rotatebox{70}{\textbf{Auto compile}} & \rotatebox{70}{\textbf{Counterex.}} & \rotatebox{70}{\textbf{CQL/FHIR}} \\
\midrule
TMR & $\times$ & $\times$ & Pair & $\times$ & $\times$ & $\times$ & $\times$ \\
MitPlan & $\checkmark$ & $\times$ & N-way & $\times$ & $\times$ & Part & $\times$ \\
GLARE & $\checkmark$ & $\times$ & N-way & $\times$ & $\times$ & $\times$ & $\times$ \\
UPPAAL & $\checkmark$ & $\times$ & Pair & $\checkmark$ & $\times$ & $\checkmark$ & $\times$ \\
Lexicomp & $\times$ & $\times$ & Pair & N/A & N/A & $\times$ & $\times$ \\
Micromedex & $\times$ & $\times$ & Pair & N/A & N/A & $\times$ & $\times$ \\
DrugBank & $\times$ & $\times$ & Pair & N/A & N/A & $\times$ & $\times$ \\
CPOE & $\times$ & $\times$ & Pair & N/A & N/A & $\times$ & Part \\
\textbf{GuardPh.} & $\checkmark$ & $\checkmark$ & N-way & $\checkmark$ & $\checkmark$ & $\checkmark$ & $\checkmark$ \\
\bottomrule
\end{tabular}
\end{table}

\textbf{Drug interaction databases} (Lexicomp~\cite{lexicomp}, Micromedex~\cite{micromedex}, DrugBank~\cite{drugbank}) maintain curated drug-drug interaction monographs. They check pairwise interactions at prescription time without temporal reasoning or formal verification. Alert fatigue from high false-positive rates (up to 90\% override rates~\cite{alert_fatigue}) is well-documented.

\textbf{CPOE systems} (Epic, Oracle Health/Cerner) integrate drug interaction databases into computerized provider order entry. Their alerts are pairwise, atemporal, and database-driven. No CPOE system performs guideline-level verification.

\textbf{TMR}~\cite{zamborlini2017} uses first-order Semantic Web rules to detect guideline interactions. TMR is atemporal and handles only pairwise interactions.

\textbf{MitPlan}~\cite{wilk2021mitplan} uses AI planning with durative operators to detect and mitigate N-guideline conflicts. It performs heuristic search, not exhaustive verification.

\textbf{GLARE/META-GLARE}~\cite{terenziani2023} executes multiple guidelines concurrently and detects conflicts via constraint propagation along a single execution path.

\textbf{Formal PK analysis.} SimBiology (MATLAB), NONMEM, and Monolix perform PK/PD simulation but are not integrated with guideline verification. dReach~\cite{gao2013delta} performs $\delta$-reachability for hybrid systems; \textsc{GuardPharma} builds on this with domain-specific PTA formalism.

% ============================================================================
\section{Threats to Validity}
\label{sec:threats}
% ============================================================================


\textbf{External validity.} Our benchmark covers only 10 scenarios derived from 6 guidelines across 4 therapeutic areas. This is a small benchmark; results may not generalize to broader polypharmacy settings (oncology, psychiatry, pediatrics). The corpus is US-focused. Larger benchmarks with clinician-adjudicated ground truth are needed.

\textbf{Benchmark authorship and circularity.} All 10 benchmark scenarios (and the 20-scenario pharmacodynamic expansion) were authored by the tool developers, not drawn from independent clinical datasets, patient records, or external pharmacovigilance databases such as FAERS. Ground truth labels were assigned by the authors rather than by independent clinical adjudicators. Achieving 100\% accuracy on a small, developer-authored test suite is expected and does not constitute clinical validation. There is an inherent risk of circular validation: scenarios may have been implicitly or explicitly selected to match the tool's detection capabilities. The separate 500-profile validator run materially reinforces this concern: despite the perfect small-benchmark scores, the saved repository output reports 319 false negatives and zero true positives against its external interaction knowledge base. Future work should evaluate on independently curated benchmarks, ideally with blinded clinical adjudication and held-out test sets not seen during development.

\textbf{No real patient data.} All scenarios use synthetic patient profiles. The benchmark does not evaluate performance on real electronic health record data, which would include noise, missing data, off-label prescribing, and patient-specific pharmacokinetic variability beyond the population-level intervals modeled here. No cross-validation or statistical confidence intervals are reported given the small sample size, and even the larger 500-profile validation run remains synthetic rather than clinical.

\textbf{Construct validity.} PK model fidelity depends on published population PK parameters representing average patients. We partially address this via interval-valued parameters covering the 5th--95th percentile range.

\textbf{CQL coverage.} Of 12 guidelines, 7 are manually encoded. We validated manual encodings against published algorithms with clinical domain experts.

% ============================================================================
\section{Discussion}
\label{sec:discussion}
% ============================================================================

\subsection{Why Temporal Reasoning Matters}







\subsection{Limitations of Contract-Based Decomposition}

While contract-based compositional verification achieves dramatic scalability improvements, it has fundamental limitations that must be acknowledged.

First, the decomposition assumes that guideline interactions are mediated primarily through shared CYP-enzyme metabolic pathways. We address some pharmacodynamic interactions---serotonin syndrome, additive CNS depression, and additive bleeding risk---via pattern-based PD pathway detection. However, other PD mechanisms (e.g., additive QT prolongation from multiple non-CYP drugs, shared renal elimination, protein binding displacement cascades) are not yet captured by the PD pathway module and must be handled by the CEGAR fallback. Extending the PD pathway catalog is straightforward engineering but requires careful domain-expert review of each rule.



To soundly handle non-Metzler interactions introduced by CYP inhibition or induction, \textsc{GuardPharma} decomposes the system Jacobian $A = M + P$, where $M$ is the Metzler part (negative off-diagonals zeroed) and $P$ is the perturbation containing the non-Metzler entries. Using the matrix measure bound, the reachable set of the perturbed system is over-approximated by:
\[
\| e^{At} x_0 - e^{Mt} x_0 \| \;\leq\; \|x_0\| \cdot t \cdot \|P\|_\infty \cdot e^{\mu_\infty(M)\,t}
\]

\subsection{Clinical Deployment Considerations}

Deploying \textsc{GuardPharma} in a clinical setting raises several practical considerations beyond the technical evaluation.

\textbf{Integration pathway.} The most natural integration point is during guideline publication and CDS configuration, not at prescription time. Hospital CDS committees would run \textsc{GuardPharma} when configuring guideline sets for their patient populations, receiving a pre-deployment safety report. This avoids adding latency to the clinical workflow and focuses the tool on its strength: exhaustive pre-deployment verification.

\textbf{Handling guideline updates.} Clinical guidelines are updated every 1--3 years. When a guideline is updated, \textsc{GuardPharma} must re-verify all interacting guideline pairs. The compositional architecture helps: if guideline $G_i$ is updated but its enzyme-interface contracts remain valid, verification of non-interacting guidelines is unaffected. Only guidelines sharing metabolic pathways with $G_i$ require re-verification.

\textbf{Population PK variability.} Our interval-valued PK parameters cover the 5th--95th percentile range of population variability. For extreme patients (e.g., neonates, morbidly obese, dialysis patients), the standard population PK models may be inadequate. Extending \textsc{GuardPharma} with specialized population PK models for these subpopulations is a natural next step.

\textbf{Regulatory pathway.} \textsc{GuardPharma}'s verification results could complement existing CDS safety standards (ONC certification criteria, Joint Commission CDS standards) by providing formal evidence of guideline compatibility. The current regulatory framework does not require formal verification of CDS interactions, but the increasing adoption of FHIR-based CDS under the 21st Century Cures Act creates a natural opportunity for verification standards.

\subsection{Comparison with LLM-Based Approaches}

Large language models with retrieval-augmented access to drug interaction databases (e.g., GPT-4 with DrugBank retrieval) can flag most known drug-drug interactions with mechanism explanations. For individual prescription checking, this approach is often practical and sufficient. However, LLMs fundamentally cannot perform the \emph{exhaustive temporal verification} that \textsc{GuardPharma} provides:

\begin{enumerate}
\item LLMs cannot prove that \emph{no} patient trajectory under a set of guidelines ever reaches a dangerous state. They can check specific scenarios but cannot enumerate all possible trajectories.
\item LLMs cannot produce \emph{sound counterexamples}---provably reachable violating traces with quantitative PK concentrations along the trajectory.
\item LLMs are stochastic: the same query may yield different results on different runs. In a safety-critical regulatory context, non-deterministic verification is inadequate.
\item LLMs cannot perform the compositional decomposition that enables scaling to 8+ guidelines; they would need to reason about all drug combinations simultaneously.
\end{enumerate}

\textsc{GuardPharma} and LLM-based approaches are complementary: LLMs excel at natural-language explanation and individual case consultation; \textsc{GuardPharma} provides the exhaustive formal guarantee layer above.

% ============================================================================
\section{Conclusion}
\label{sec:conclusion}
% ============================================================================


The key technical enablers are the PTA formalism (integrating pharmacokinetic dynamics into formal verification), enzyme-interface contracts (exploiting metabolic pathway structure for compositional reasoning), and the $\delta$-decidability result (providing sound reachability analysis at clinically relevant precision). In the current repository state, these ideas look strongest as a research prototype: they perform well on small developer-authored benchmarks and the pharmacodynamic extension, but the broader 500-profile validator run shows that substantial work remains before the system can be described as clinically reliable.

\subsection*{Data Availability}
The evaluation scripts and saved result files discussed in this paper are included in this repository under \texttt{benchmarks/} and \texttt{benchmarks/results/}. All patient scenarios use synthetic data.

% ============================================================================
\begin{thebibliography}{99}

\bibitem{cms2019chronic}
Centers for Medicare and Medicaid Services.
Chronic conditions among Medicare beneficiaries, 2019.
CMS Chartbook, 2019.

\bibitem{budnitz2011emergency}
D.~S. Budnitz, M.~C. Lovegrove, N.~Shehab, and C.~L. Richards.
Emergency hospitalizations for adverse drug events in older Americans.
\emph{New England Journal of Medicine}, 365(21):2002--2012, 2011.

\bibitem{boyd2005clinical}
C.~M. Boyd et al.
Clinical practice guidelines and quality of care for older patients with multiple comorbid diseases.
\emph{JAMA}, 294(6):716--724, 2005.

\bibitem{qato2016changes}
D.~M. Qato et al.
Changes in prescription and over-the-counter medication and dietary supplement use among older adults in the United States.
\emph{JAMA Internal Medicine}, 176(4):473--482, 2016.

\bibitem{zamborlini2017}
V.~Zamborlini et al.
Inferring recommendation interactions in clinical guidelines.
\emph{Semantic Web}, 8(3):421--446, 2017.

\bibitem{wilk2021mitplan}
S.~Wilk et al.
Comprehensive mitigation framework for concurrent application of multiple clinical practice guidelines.
\emph{Journal of Biomedical Informatics}, 113:103621, 2021.

\bibitem{terenziani2023}
P.~Terenziani et al.
Managing the interaction between multiple clinical guidelines.
\emph{Applied Sciences}, 13(4):2414, 2023.

\bibitem{uppaal_medical}
A.~David et al.
Uppaal SMC tutorial.
\emph{Int.\ Journal on Software Tools for Technology Transfer}, 17(4):397--415, 2015.

\bibitem{cql_spec}
HL7 International.
Clinical Quality Language (CQL) Specification, Release 1.5, 2023.

\bibitem{fhir_plandefinition}
HL7 International.
FHIR R4 PlanDefinition Resource, 2023.

\bibitem{rowland2011clinical}
M.~Rowland and T.~N. Tozer.
\emph{Clinical Pharmacokinetics and Pharmacodynamics}.
Lippincott Williams and Wilkins, 4th edition, 2011.

\bibitem{flockhart_cyp}
D.~A. Flockhart.
Drug interactions: cytochrome P450 drug interaction table.
Indiana University School of Medicine, 2007.

\bibitem{alur1994theory}
R.~Alur and D.~L. Dill.
A theory of timed automata.
\emph{Theoretical Computer Science}, 126(2):183--235, 1994.

\bibitem{henzinger1996theory}
T.~A. Henzinger.
The theory of hybrid automata.
In \emph{Proc.\ LICS}, pages 278--292. IEEE, 1996.

\bibitem{gao2013delta}
S.~Gao, S.~Kong, and E.~M. Clarke.
dReal: an SMT solver for nonlinear theories over the reals.
In \emph{CADE}, pages 208--214. Springer, 2013.

\bibitem{drugbank}
D.~S. Wishart et al.
DrugBank 5.0: a major update to the DrugBank database for 2018.
\emph{Nucleic Acids Research}, 46(D1):D1074--D1082, 2018.

\bibitem{lexicomp}
Wolters Kluwer.
Lexicomp Online, 2024.

\bibitem{micromedex}
IBM Watson Health.
Micromedex Solutions, 2024.

\bibitem{beers2023}
AGS 2023 Beers Criteria Update Expert Panel.
AGS 2023 updated Beers Criteria for potentially inappropriate medication use in older adults.
\emph{J.\ Am.\ Geriatrics Soc.}, 71(7):2052--2081, 2023.

\bibitem{alert_fatigue}
H.~van der Sijs et al.
Overriding of drug safety alerts in computerized physician order entry.
\emph{JAMIA}, 13(2):138--147, 2006.

\bibitem{uw_didb}
University of Washington.
Drug Interaction Database, 2024.

\bibitem{dumbreck2015drug}
S.~Dumbreck et al.
Drug-disease and drug-drug interactions in 12 UK national clinical guidelines.
\emph{BMJ}, 350:h949, 2015.

\bibitem{kamke1932}
E.~Kamke.
Zur Theorie der Systeme gewoehnlicher Differentialgleichungen II.
\emph{Acta Mathematica}, 58:57--85, 1932.

\bibitem{smith2008monotone}
H.~L. Smith.
\emph{Monotone Dynamical Systems}.
AMS, 2008.

\bibitem{faers_data}
FDA.
FDA Adverse Event Reporting System (FAERS), 2024.

\bibitem{cqframework}
CQFramework.
Clinical Quality Framework Reference Implementation, 2024.

\bibitem{hapi_fhir}
Smile CDR.
HAPI FHIR, 2024.

\bibitem{z3}
L.~de Moura and N.~Bjorner.
Z3: An efficient SMT solver.
In \emph{TACAS}, pages 337--340. Springer, 2008.

\bibitem{capd}
T.~Kapela et al.
CAPD::DynSys: a flexible C++ toolbox for rigorous numerical analysis of dynamical systems.
\emph{Commun.\ Nonlinear Sci.\ Numer.\ Simul.}, 101:105578, 2021.

\end{thebibliography}

\end{document}
